\section{二点間の変化と一点での変化}\label{節:導入変化}
関数 $y=f(x)$ において，$x$ が $a$ から $a+h$ まで変わったとする。$h\ne0$ のとき，横の変化量は $h$，縦の変化量は $f(a+h)-f(a)$ だから，二点を結ぶ直線の傾きは
\[
 \frac{f(a+h)-f(a)}{h}
\]
である。これをその区間の\textbf{平均変化率}という。「平均」と呼ぶのは，途中の変化のしかたを区別せず，始点と終点の変化だけで一つの数を作っているからである。

$f(x)=x^2$，$a=1$ では，
\[
 \frac{(1+h)^2-1}{h}
 =\frac{2h+h^2}{h}=2+h \qquad(h\ne0)
\]
となる。$h=1$ なら傾きは $3$，$h=0.1$ なら $2.1$，$h=-0.1$ なら $1.9$ である。$h$ を正の側からも負の側からも $0$ に近づけると，傾きは $2$ に近づく。この値を，放物線上の点 $(1,1)$ における接線の傾きとして採用したい。

\begin{figure}[htbp]
\centering
\begin{tikzpicture}[x=2.6cm,y=1.5cm,>=Stealth]
 \draw[->] (-0.12,0)--(2.1,0) node[right]{$x$};
 \draw[->] (0,-0.25)--(0,3.75) node[above]{$y$};
 \draw[very thick,domain=0:1.9,samples=80] plot (\x,{\x*\x}) node[right]{$y=x^2$};
 \draw[blue,thick,domain=0.45:1.75] plot (\x,{2*\x-1});
 \draw[red,thick,domain=0.55:1.8] plot (\x,{2.6*\x-1.6});
 \fill (1,1) circle(1.6pt) node[below right]{$P(1,1)$};
 \fill (1.6,2.56) circle(1.6pt) node[right]{$Q$};
 \draw[dashed] (1,0) node[below]{$1$}--(1,1);
 \draw[dashed] (1.6,0) node[below]{$1+h$}--(1.6,2.56);
 \node[blue] at (1.82,2.08) {接線};
 \node[red] at (0.65,1.6) {割線};
\end{tikzpicture}
\caption{二点を結ぶ割線の傾きを，点 $Q$ を点 $P$ に近づけて調べる。}
\label{図:導入接線}
\end{figure}

図\ref{図:導入接線}のように，曲線上の二点を結ぶ直線を\textbf{割線}と呼ぶ。割線は二点で決まるのに対し，接線の傾きは一つの点の近くの関数の振る舞いで決まる。図だけを見て「曲線に一度だけ触れる直線」と定義すると，接線が別の場所で再び曲線と交わる場合を扱えない。微分では，割線の傾きの極限を使って接線を定める。

\begin{example}{平均の速さから瞬間の速さへ}
一直線上を動く物体の時刻 $t$ 秒における位置が $s(t)=t^2$ メートルで表されている。時刻 $2$ 秒から $2+h$ 秒までの平均の速度を求め，時刻 $2$ 秒における速度を予想せよ。
\end{example}
位置の変化を経過時間で割ればよいので，平均の速度は
\[
 \frac{s(2+h)-s(2)}{h}
 =\frac{(2+h)^2-4}{h}=4+h
\]
メートル毎秒である。たとえば $h=0.1$ なら $4.1$，$h=0.01$ なら $4.01$ となる。時刻 $2$ 秒より前の区間を調べるために $h<0$ としても，同じ式が使える。どちらから近づけても $4$ に近づくので，瞬間の速度は $4$ メートル毎秒と予想できる。

この計算では $h$ を $0$ に代入していない。元の式では $h=0$ にすると分母も分子も $0$ となり，割り算ができない。$h\ne0$ のまま式を整理し，$h$ が $0$ に近づくときの値を調べたのである。\textbf{変化量を $0$ にすることと，$0$ に近づけることは別の操作}である。

